\subsection*{Classical stylized facts and the failure of Gaussian diffusion}
The quantitative study of financial market dynamics has been shaped by a tension between tractable theoretical frameworks and the empirical complexity of observed data. The Black-Scholes-Merton option pricing model \cite{black_scholes_1973} and the geometric Brownian motion (GBM) paradigm embedded within it assume that log-returns are independently and identically distributed Gaussian random variables. Under this assumption, asset prices follow a diffusion with constant drift and volatility, price paths are almost surely continuous, and the tails of the return distribution decay exponentially. These properties have well-known practical advantages: GBM admits closed-form option prices, portfolio variances aggregate linearly, and model parameters can be estimated with minimal data.

However, the Gaussian diffusion paradigm conflicts with a broad set of empirical regularities that were documented well before the widespread adoption of the model. Mandelbrot \cite{mandelbrot_variation_1963} demonstrated that the distributions of speculative price changes exhibit excess kurtosis; the probability of large deviations from the mean is far higher than Gaussian predictions, a phenomenon now described as \textit{fat tails} or \textit{leptokurtosis}. Fama \cite{fama_efficient_1970} showed that, consistent with the efficient markets hypothesis, raw return series display negligible linear autocorrelation; daily or weekly returns are essentially unpredictable from their own past values. Yet Schwert \cite{schwert_why_1989} and others documented that the \textit{magnitude} of returns, measured by absolute or squared returns, exhibits substantial positive autocorrelation that decays slowly over time. This temporal clustering of variance, where large price changes tend to be followed by large changes regardless of sign, is referred to as \textit{volatility clustering} and is among the most persistently observed features of financial time series. Cont \cite{cont_empirical_2001} provides a systematic review of these and related stylized facts, confirming their presence across asset classes, geographic markets, and sampling frequencies. Collectively, leptokurtosis, the absence of return predictability, and volatility clustering define the benchmark against which any proposed model of market dynamics must be evaluated.

\subsection*{Conditional variance models: ARCH and GARCH families}
The most influential class of models designed to address volatility clustering within the Gaussian framework is the autoregressive conditional heteroscedasticity (ARCH) family introduced by Engle \cite{engle_1982}. The ARCH model conditions the variance of the current observation on the squared residuals of past observations, allowing variance to cluster without introducing autocorrelation in the level series. Bollerslev \cite{bollerslev_garch_1986} generalized this to the GARCH(p,q) framework, in which the conditional variance depends on both past squared residuals and past conditional variances; the GARCH(1,1) specification became the canonical one-factor volatility model and remains widely used in practice. GARCH models are parsimonious, analytically tractable, and can be estimated by maximum likelihood under distributional assumptions on the innovation.

Despite their success in modeling conditional variance dynamics, standard GARCH models have several well-documented limitations in the context of stylized-facts reproduction. First, while GARCH innovations can be made leptokurtic by specifying Student-t or generalized error distributions, the fat-tail behavior arises through the parametric distributional assumption rather than through the structural dynamics of the model. Second, GARCH models do not explicitly represent discrete market regimes or the possibility of sudden jumps in the volatility process; the variance responds smoothly and continuously to past shocks. Third, regime persistence, the tendency of markets to remain in high-volatility states for sustained periods, is captured only indirectly through the persistence parameter of GARCH, which governs exponential decay of the volatility response. When persistence is near unity, the GARCH model is near-integrated and exhibits long-memory-like behavior, but this is a parametric approximation rather than a structural representation of distinct market regimes \cite{diebold_inoue_2001}. Stochastic volatility models \cite{stein_stock_1991, heston_stochastic_1993} address some of these shortcomings by treating variance as a latent diffusion process, but they require computationally intensive estimation via particle filters or Markov chain Monte Carlo methods, and their parameters do not directly correspond to observable market states.

\subsection*{Jump-diffusion models}
An alternative approach to modeling fat tails and discontinuous price behavior is the explicit introduction of Poisson-driven jumps into the asset price process. Merton \cite{merton_jump_1976} extended the Black-Scholes framework by adding a compound Poisson jump component to the GBM, producing a model in which prices follow a diffusion between events and undergo discrete jumps at random times governed by a Poisson arrival process. This construction directly addresses leptokurtosis: the marginal distribution of returns is a Gaussian mixture weighted by the Poisson probability of observing zero, one, or more jumps over the sampling interval. Option pricing under Merton's model requires integrating over the jump count distribution, but retains analytical tractability through the characteristic function.

Kou \cite{kou_jump-diffusion_2002} extended the jump-diffusion framework by replacing the Gaussian jump amplitude distribution with a double-exponential (asymmetric Laplace) distribution, which provides a better fit to the empirically observed asymmetry between the sizes of positive and negative jumps. This modification preserves analytical tractability for option pricing while improving the empirical fit to return distributions. Jump-diffusion models are particularly relevant in the context of equity markets, where identifiable macroeconomic announcements, earnings surprises, and geopolitical events generate discrete, large-magnitude price movements that are distinct from the continuous diffusion component \cite{au_yeung_jump_2020}. A fundamental challenge, however, is that jump-diffusion models in their standard form do not provide an explicit mechanism for volatility persistence; the post-jump volatility immediately returns to its baseline level, absent a separate stochastic volatility component. Integrating jump-diffusion with regime-switching provides a natural solution to this limitation, motivating the hybrid framework developed in the present study.

\subsection*{Hidden Markov models in finance}
Hidden Markov models provide a probabilistic framework for non-stationary time series in which the observed data are generated by a process that switches among a discrete set of latent states according to Markov dynamics. The foundational algorithmic framework was developed by Rabiner and Juang \cite{rabiner_introduction_1986}, who introduced the Baum-Welch algorithm for maximum-likelihood estimation of the emission and transition parameters from observed sequences. Hamilton \cite{hamilton_1989} applied the regime-switching HMM to macroeconomic time series, demonstrating that U.S. GDP growth is well described by a two-state model with distinct expansion and recession regimes; this work established the template for applying latent-state models to financial time series. The identification of time-varying market regimes has since become a central application of HMMs in finance, with applications ranging from stock return modeling to volatility forecasting and asset allocation.

Rydén, Teräsvirta, and \AA sbrink \cite{ryden_stylized_1998} provided an influential systematic evaluation of HMM specifications against the stylized facts of daily return series. They showed that simple Gaussian-emission HMMs with a small number of states could reproduce leptokurtosis and the absence of return autocorrelation, but failed to generate the observed volatility clustering; ACF of absolute returns under standard HMMs decayed far too rapidly relative to empirical data. This diagnostic limitation identified the core structural gap that the present study addresses through the jump-duration extension. Bulla and Bulla \cite{bulla_stylized_2006} advanced this line of inquiry by considering hidden semi-Markov models, in which the dwell time in each state follows a state-specific distribution rather than the geometric distribution implied by standard Markov dynamics. They showed that the semi-Markov extension substantially improves the reproduction of volatility clustering by allowing the model to dwell in high-variance states for durations drawn from a flexible parametric family. Rossi and Gallo \cite{rossi_volatility_2006} applied HMMs to multivariate volatility estimation, demonstrating that latent state dynamics can capture the joint clustering of volatility across assets in a more interpretable manner than multivariate GARCH models. Nguyen \cite{nguyen_hidden_2018} further demonstrated the utility of HMMs for stock trading, showing that regime identification from price data can inform tactical asset allocation.

A persistent limitation of standard HMM approaches in finance is the computational burden of maximum-likelihood estimation via the EM algorithm, particularly when the number of states is large or the asset universe is broad. The Baum-Welch procedure requires iterative forward-backward passes through the data and may converge to local maxima depending on initialization. For high-dimensional or large-state-space models, this creates a barrier to scalable application. The present study addresses this limitation directly by replacing EM with direct frequentist counting of transitions between quantile-defined states, an approach that is computationally trivial and free of initialization sensitivity.

\subsection*{Synthetic data generation and generative frameworks}
The generation of synthetic financial time series has attracted increasing attention as both a practical tool (for augmenting limited data, testing risk models, and stress testing portfolios) and as a benchmark for evaluating the quality of generative models. Assefa et al.\ \cite{assefa_generating_2021} provide a comprehensive survey of the opportunities, challenges, and pitfalls of synthetic data generation in finance, noting that stylized-facts reproduction is the primary criterion for evaluating the statistical quality of synthetic financial time series, and that naive generative approaches tend to fail on at least one of the three core stylized facts.

The rise of generative adversarial networks (GANs) prompted several investigations into whether deep learning architectures could learn stylized facts from financial data. Takahashi, Chen, and Tanaka-Ishii \cite{takahashi_2019} showed that standard GAN architectures produce synthetic return sequences that visually resemble financial time series but fail to reproduce ACF structure precisely. Kwon and Lee \cite{kwon_can_2024} provided a careful evaluation of GAN-based methods against the stylized-facts checklist, finding that while GANs match marginal distributions reasonably well, capturing temporal dependence structures, particularly volatility clustering, remains challenging without explicitly encoding the relevant inductive biases into the architecture or training objective. Foundation models for Markov jump processes, as introduced by Berghaus et al.\ \cite{berghaus_foundation_2025}, offer a complementary perspective: large-scale pretrained inference networks can be applied to new Markov process instances without retraining, but their application to financial time series generation remains an open research direction.

Against this backdrop, the hybrid HMM framework of the present study occupies a distinct position: it is explicitly designed around the structural properties known to generate stylized facts, uses an interpretable discrete-state representation that can be audited and stress-tested, and achieves strong statistical fidelity with minimal computational overhead. The empirical estimation procedure avoids the training instability and data requirements of deep generative models, while the jump-duration mechanism provides a principled solution to the volatility-clustering limitation of standard HMMs. These properties make the framework particularly well-suited for deployment in risk management and stress-testing pipelines, where interpretability and computational reproducibility are as important as statistical fidelity.

\subsection*{Factor models and multi-asset extensions}
Modeling the joint dynamics of large asset universes requires a framework for managing cross-sectional dependence. The Single-Index Model (SIM) of Sharpe \cite{sharpe_simplified_1963} provides the canonical linear factor decomposition: each asset's excess return is decomposed into a common market factor component and an idiosyncratic residual, drastically reducing the number of parameters required to specify a covariance matrix relative to a full multivariate model. The SIM remains foundational in portfolio construction and risk management despite the availability of richer multi-factor models such as Fama-French \cite{fama_french_1993}, because its simplicity enables transparent attribution and scalable computation.

Existing HMM-based work in finance has been predominantly univariate, focusing on single-asset or single-index regime identification. Applications to individual equities \cite{nguyen_hidden_2018}, volatility modeling \cite{rossi_volatility_2006}, and macro time series \cite{hamilton_1989, kim_nelson_1999} collectively demonstrate the utility of latent state representations, but leave open the question of how HMM-derived generative models can be extended to large cross-sections of assets in a computationally tractable way. Mixture hidden Markov models \cite{dias_mixture_2010} offer a partial solution by allowing multiple assets to share a common latent state structure, but estimating such models jointly across hundreds of assets is computationally demanding.

The present study addresses this gap by combining the univariate HMM framework with the SIM factor decomposition. The hybrid HMM is estimated on the SPY index, which serves as the market-factor generator; asset-level paths are then reconstructed via the linear SIM projection. This construction exploits the fact that the dominant source of joint volatility clustering in equity markets is the common factor, and that idiosyncratic shocks can be treated as approximately independent given the factor. The result is a scalable generative pipeline for a 424-asset universe that inherits the stylized-facts fidelity of the HMM without requiring joint estimation of a high-dimensional model.
